% ==============================================================================
% THIETKE.TEX - HỆ THỐNG THIẾT KẾ & ĐỊNH DẠNG ĐỀ THI
% Hỗ trợ 3 chế độ: 1 = Giáo viên | 2 = HS dòng chấm | 3 = HS chỉ in đề
% ==============================================================================

% ============================== 1. QUẢN LÝ CHẾ ĐỘ ==============================
\newif\ifHS
\newif\ifDongCham
\newif\ifDapAn

\ifnum\chedo=1
	\HSfalse \DongChamfalse \DapAntrue     % Chế độ Giáo viên (đầy đủ lời giải + đáp án)
\else\ifnum\chedo=2
	\HStrue  \DongChamtrue  \DapAntrue    % Chế độ HS dòng chấm (kẻ dòng chấm cho lời giải)
\else
	\HStrue  \DongChamfalse \DapAnfalse    % Chế độ HS chỉ đề (gọn gàng, không lời giải)
\fi\fi

% ============================== 2. GÓI TOÁN HỌC & FONT =========================
\usepackage{amsmath,amssymb,yhmath,mathrsfs}
\usepackage{unicode-math}
\usepackage{etoolbox}
\usepackage{esvect}
% unicode-math ghi đè \overrightarrow lúc \begin{document}, nên đặt lại bằng \AtBeginDocument
\AtBeginDocument{\def\overrightarrow{\vv}}

% Font & Ngôn ngữ (XeLaTeX)
\usepackage{polyglossia}
\setmainlanguage{vietnamese}
\setmainfont{TeX Gyre Termes}
\setsansfont{Arial}
\setmonofont{Consolas}
\setmathfont{Cambria Math}

% ============================== 3. ĐỒ HỌA, BẢNG BIỂU & TRANG ===================
\usepackage{graphicx,xcolor}
\usepackage{wrapfig,varwidth,diagbox,multirow,makecell}
\usepackage{tikz,tikz-3dplot,tkz-tab,tkz-euclide}
\usetikzlibrary{math,shadows,through,fadings,shadings,quotes,shapes,decorations.text,matrix,calc,angles,arrows.meta,intersections,patterns}

\usepackage{pifont,bbding,fontawesome,pgfornament}
\usepackage{enumerate}
\usepackage[top=1.5cm,bottom=1.5cm,left=1.4cm,right=1.4cm]{geometry}
\usepackage{fancyhdr,lastpage}
\usepackage[hidelinks,unicode,hypertexnames=false]{hyperref}

% ============================== 4. TRẮC NGHIỆM & DÒNG CHẤM ====================
\ifHS\usepackage[dethi]{ex_test}\else\usepackage[solcolor]{ex_test}\fi
\usepackage{framed}
\usepackage[many]{tcolorbox}
\usepackage[outline]{contour}
\usepackage{ex_dotfit}            % nạp SAU ex_test

% Khóa câu hỏi trả lời ngắn không hiện khung ở chế độ 3
\ifnum\chedo=3 \OPTN{kindSA=0}\fi

% ----- Bản vá tương thích XeTeX & SyncTeX của ex_test -----
\makeatletter
% Ngăn ex_test tạm tắt SyncTeX (\synctex=0) trong lệnh \choice
% Tránh lỗi XeTeX vô hiệu hóa vĩnh viễn SyncTeX khi ngắt trang ở chế độ 3
\def\extest@choiceSyncMuteBeginActive{}
\def\extest@choiceSyncMuteEndActive{}
\makeatother

\renewcommand{\textchoice}{%
	\par\vspace{2pt}\noindent%
	\ifnum\chType=0\relax\selectchoice\ \keyCh%
	\else\ifcase\numexpr\chType+3\relax
	\selectchoiceTF\ \keyTF \or
	\selectchoiceTF\ \keyTF \or
	\selectshortans\ \keySA \or\or
	\selectchoiceTF\ \keyTF \or
	\selectchoice\ \keyTF
	\fi\fi%
	\finishLG%
}

% ----- Cấu hình dòng chấm (tự động kích hoạt khi ở Chế độ 2) -----
\newcommand{\dotscale}{0.6}
\newcommand{\dotmin}{2}
\newcommand{\dotmax}{6}

\ifDongCham
	\dotplanTF{2}{3}
	\dotplanSA{2}{4}
	\dotplanmax{14}
	\AtBeginDocument{%
		\dotlinefull{ex}[scale=\dotscale,minlines=\dotmin,maxlines=\dotmax]%
	}
\fi

% ============================== 5. MÀU SẮC & GIAO DIỆN =========================
\definecolor{mauchinh}{RGB}{45,81,133}
\definecolor{xanhlo}{rgb}{0.63,0.84,0.71}
\definecolor{doruby}{rgb}{1.0,0.0,0.16}
\definecolor{xanhladam}{rgb}{0.0,0.26,0.26}

\linespread{1.25}
\newenvironment{tomtat}{\begin{tcolorbox}[boxrule=1pt,arc=0mm,breakable]}{\end{tcolorbox}}
\newcommand{\hoac}[1]{\left[\begin{aligned}#1\end{aligned}\right.}
\newcommand{\heva}[1]{\left\{\begin{aligned}#1\end{aligned}\right.}
\def\qedEX{}
\newtheorem{vd}{Ví dụ}
\newtheorem{bt}{Bài}

% Định dạng \subsection* (khung bo tiêu đề phần chạy cùng dòng)
\usepackage{titlesec}
\newtcbox{\khungtieude}{on line,colback=white,colframe=black,boxrule=0.8pt,
	sharp corners,boxsep=0pt,left=4pt,right=4pt,top=3pt,bottom=3pt}
\titleformat{\subsection}[runin]{\normalfont\bfseries}{}{0pt}{\khungtieude}
\titlespacing*{\subsection}{0pt}{10pt}{0.6em}

% ============================== 6. CHÂN TRANG & MÃ ĐỀ ==========================
\def\madethi{}                              % Do \nhapde tự reset; file đề dùng \MaDe{...}
\newcommand{\MaDe}[1]{\gdef\madethi{#1}}

% Đếm số thứ tự đề để hiển thị đúng tổng số trang riêng của đề đó
\newcounter{sode}
\pagestyle{fancy}
\fancyhf{}
\renewcommand{\headrulewidth}{0pt}
\renewcommand{\footrulewidth}{0pt}
\fancyfoot[R]{Trang \thepage/\pageref{fin-\thesode}%
	\ifx\madethi\empty\else\ - Mã đề thi \madethi\fi}

% ============================== 7. QUẢN LÝ NHẬP ĐỀ THI & ĐÁP ÁN ==============
\def\thumucde{}                             % Thư mục chứa file đề, vd: data/HSG12/
\newcommand{\cotPI}{12}
\newcommand{\cotPII}{6}
\newcommand{\cotPIII}{6}

% Tít đề 2 cột
\newcommand{\TitDe}[2]{%
	\begin{center}
		\begin{minipage}[t]{0.45\textwidth}\centering #1\end{minipage}\hfill
		\begin{minipage}[t]{0.55\textwidth}\centering #2\end{minipage}
	\end{center}
	\ifx\madethi\empty\else
	\par\vspace{-2pt}\hfill\fbox{\textbf{\ Mã đề thi: \madethi\ }}\par\vspace{4pt}
	\fi
}

\makeatletter
\newwrite\xoa@w
% Xóa file đáp án tạm của đề trước (đảm bảo đề thiếu phần không bị đè đáp án cũ)
\def\xoaans#1{\immediate\openout\xoa@w=ans/ans-#1.tex \immediate\closeout\xoa@w}

% In một phần đáp án CHỈ KHI file đáp án có ít nhất một câu
\newcommand{\inphan}[3]{%
	\extest@probe{ans/ans-#2}%
	\ifnum\extest@seen=1 \noindent\textbf{#1}\inputanstab{#3}{ans/ans-#2}\fi}

% Lệnh nhập 1 file đề
\newcommand{\nhapde}[1]{%
	\stepcounter{sode}%
	\setcounter{page}{1}%
	\gdef\madethi{}%
	\xoaans1\xoaans2\xoaans3\xoaans4%
	\input{\thumucde #1}%
	\clearpage
	\immediate\write\@auxout{\string\newlabel{fin-\arabic{sode}}{{}{\number\numexpr\value{page}-1\relax}{}{}{}}}%
}
\makeatother

% Nhập danh sách nhiều file đề: \nhapdanh{de1,de2,de3}
\newcommand{\nhapdanh}[1]{\forcsvlist{\nhapde}{#1}}

% In đáp án ở cuối mỗi file đề (được gọi từ file đề)
\newcommand{\hethet}{\begin{center}-----------------------HẾT-----------------------\end{center}}
\newcommand{\dapan}{%
	\hethet
	\ifDapAn
	%\newpage
	\begin{center}\textbf{ĐÁP ÁN}\end{center}
	\rowans
	\inphan{PHẦN I}{1}{\cotPI}
	\inphan{PHẦN II}{2}{\cotPII}
	\inphan{PHẦN III}{3}{\cotPIII}
	\fi
}